\section{Udspænding af vektorer}
En udspænding af vektorer er mængden af alle vektorer der kan dannes ved linearkombinationer af et givent sæt vektorer. (hvor vektorer ganges med $c \in \mathbb{F}$ og lægges sammen)

Hvis en af vektorerne kan skrives some en linearkombination af de andre vektorer, er vektorene lineært afhængige, og denne vektor kan fjernes uden at ændre udspændingen. Ellers er vektorene lineært uafhængige.

\subsection{Ordnet basis for udspænding af vektorer}
Alle vektorer i en ordnet basis er lineært uafhængige. 
\kilde{Korollar 9.2.2}

\label{sec:ordnet-basis-udspaending}
\subsubsection{Find ordnet basis for udspænding}
Lav en matrix ud af vektorene

\[
    \begin{bmatrix}
        \vline & \vline & \cdots & \vline \\
        \vec{v_1} & \vec{v_2} & \cdots &  \vec{v_n} \\
        \vline & \vline & \cdots & \vline \\

    \end{bmatrix}
\]

Brug rækkeoperationer til at få matrixen på trappeform. 
For alle søjler med et pivot-element skal den \emph{originale vektor} på den position med i den ordnede basis.

\kilde{Sætning 9.2.1}

\subsection{Er vektor i udspænding?}
\label{sec:span}

\subsubsection{Determinant}
Hvis determinanten af matrixen skabt ved at tilføje den nye vektor som en ekstra søjle er 0, så er vektoren indeholdt i spandet.


\subsubsection{Rækkeoperationer}
Lav en matrix med vektorerne i spanet og den nye vektor som ekstra søjle:
\[
    \begin{bmatrix}
        \vline & \vline & \cdots & \vline & \vline \\
        \vec{v_1} & \vec{v_2} & \cdots & \vec{v_n} & \vec{v_x} \\
        \vline & \vline & \cdots & \vline & \vline \\
    \end{bmatrix}
\]
Få matricen på reduceret rækkeform. Hvis den ekstra søjle \emph{ikke} indeholder et pivot-element, er vektoren indeholdt i udspændingen.

I så fald kan vektorens linearkombinationen aflæses fra søjlen $\vec{v_x}$


\eksempel{
\[
\begin{bmatrix}
    1 & 3 & 7 \\
    2 & 3 & 8 \\
    1 & 0 & 1 \\
    4 & 1 & 6
\end{bmatrix}
\xrightarrow{rref}
\begin{bmatrix}
    1 & 0 & 1 \\
    0 & 1 & 2 \\
    0 & 0 & 0 \\
    0 & 0 & 0
\end{bmatrix}
\]

Her kan aflæses at \(
    \vec{v_x} = 1 \cdot \vec{v_1} + 2 \cdot \vec{v_2}
\)
}

\subsection{Vektor med ubekendt i udspænding}
Lav en matrix med vektorerne og den nye vektor som ekstra søjle:
\[
    \begin{bmatrix}
        \vline & \vline & \cdots & \vline & \vline \\
        \vec{v_1} & \vec{v_2} & \cdots & \vec{v_n} & \vec{v_x} \\
        \vline & \vline & \cdots & \vline & \vline \\
    \end{bmatrix}
\]
Brug rækkeoperationer til at få matrixen på trappeform. Løs for den ubekendte så den ekstra søjle \emph{ikke} indeholder et pivot-element, så er vektoren indeholdt i udspændingen.

\subsection{Dimensionssætningen}
For en matrix \( \mathbf{A}^{m \times n} \) ($n$ søjler) gælder det at:
\[
    \rho(\mathbf{A}) + \dim(\ker(\mathbf{A})) = n \kildetag{Sætning 9.4.2}
\]
$\rho(\mathbf{A}) = \dim(\colsp(\mathbf{A})) =$ antal pivoter \hfill $\dim(\ker(\mathbf{A})) = \nullop(\mathbf{A})$

\

Dimensionen er antallet af vektorer i en ordnet basis for rummet.

\subsubsection{Difference mellem vektorer}
Hvis \(\mathbf{A} \cdot \vec{v_1} = \mathbf{A} \cdot \vec{v_2}\) er \(c \cdot (\vec{v_1}-\vec{v_2})\) i \(\ker(\mathbf{A})\) 

\subsection{Søjle-rum}
\label{sec:soejlerum}

Mængden af alle linearkombinationer af søjlerne i en matrix.

Spanet over alle søjler med et pivot-element.

\[
    \colsp(\mathbf{A}) = \{ \mathbf{A} \cdot \vec{v} \ | \ \vec{v} \in \mathbb{F}^n \}
\]

Ordnet basis for søjlerummet kan findes som beskrevet i sektion \ref{sec:ordnet-basis-udspaending}.

\subsection{Kernen / Nulrummet af en matrix}
\label{sec:kerne}

\[ \ker(\mathbf{A}) = \{ \vec{v} \in \mathbb{F} \ | \ \mathbf{A} \cdot \vec{v} = 0 \} \]

Kernen for matrix \textbf{A} findes ved at løse det Hom.LLS dannet ved at tilføje en 0-søjle til højre.
Se sektion \ref{sec:LLS-løsning} for løsning af LLS.

\kilde{Korollar 9.2.3}

%\eksempel{
%    \[
%        \begin{bmatrix}
%            -i & 1 & 0 & \vline & 0 \\
%            2 & -2 & 2 & \vline & 0 \\
%            i & -1 & 0 & \vline & 0
%        \end{bmatrix}
%    \]
%
%    Første linje giver os
%    \[
%        -i \cdot x_1 + x_2 = 0 \implies x_2 = i \cdot x_1
%    \]
%
%    Nu kan vi finde $x_3$ med anden linje, hvor vi udskifter $x_2$ med $i \cdot x_1$:
%    \begin{align*}
%        2 \cdot x_1 - 2 \cdot (i \cdot x_1) + 2 \cdot x_3 &= 0 \\
%        \implies 2 \cdot x_1 - 2i \cdot x_1 &= -2 \cdot x_3 \\
%        \implies x_1 - i \cdot x_1 &= -x_3 \\
%        \implies (1 - i) \cdot x_1 &= -x_3 \\
%        \implies (i - 1) \cdot x_1 &= x_3
%    \end{align*}
%
%    Vi har nu både $x_2$ og $x_3$ udtrykt ved $x_1$, så vi kan skrive løsningen som:
%
%    \[
%        \begin{bmatrix}
%            x_1 \\
%            x_2 \\
%            x_3
%        \end{bmatrix}
%        =
%        \begin{bmatrix}
%            x_1 \\
%            i \cdot x_1 \\
%            (i - 1) \cdot x_1
%        \end{bmatrix}
%        =
%        x_1 \cdot
%        \begin{bmatrix}
%            1 \\
%            i \\
%            i - 1
%        \end{bmatrix}
%    \]
%
%    Kernen er derfor
%    \[
%        \spanop
%        \left(
%            \begin{bmatrix}
%                1 \\
%                i \\
%                i - 1
%            \end{bmatrix}
%        \right)
%    \]
%}

\subsubsection{Ordnet basis for kernen / nulrummet}
En ordnet basis for kernen findes ved spandet over de parametre-afhlængige vektorer i løsningen til det tilsvarende Hom.LLS.

Givet vektorer i uordnet kerne kan ordnet basis findes som i sektion \ref{sec:ordnet-basis-udspaending}.

\subsubsection{Er vektor i kernen af en matrix?}
$\vec{v}$ er i kernen af matrix $\mathbf{A}$ hvis:
\begin{itemize}
    \setlength{\itemsep}{0em}
    \item $\mathbf{A} \cdot \vec{v} = 0$
    \item $\vec{v}$ er i udspændingen af vektorerne i kernen (Se sektion \ref{sec:span})
\end{itemize}


\subsection{Koordinater med hensyn til ordnet basis}
Lav en matrix med basisvektorerne og vektoren der skal findes koordinater for:
\[
    \begin{bmatrix}
        \vline & \vline & \cdots & \vline & \vline \\
        \vec{b_1} & \vec{b_2} & \cdots & \vec{b_n} & \vec{v} \\
        \vline & \vline & \cdots & \vline & \vline \\
    \end{bmatrix}
\]

Omskriv til \emph{rref}, sidste søjle er koordinaterne med hensyn til basisen.

\eksempel{

    Givet \( \vec{v} = \begin{pmatrix}
        5 \\
        7 \\
        4
    \end{pmatrix} \), \quad \( \beta = \left\{ \begin{pmatrix}
        1 \\
        1 \\
        0
    \end{pmatrix}, \begin{pmatrix}
        0 \\
        1 \\
        1
    \end{pmatrix}, \begin{pmatrix}
        1 \\
        0 \\
        1
    \end{pmatrix} \right\} \), findes \( \vec{[v]_\beta} \) ved at løse LLS'et:
    \[
        \begin{bmatrix}
            1 & 0 & 1 & \vline & 5 \\
            1 & 1 & 0 & \vline & 7 \\
            0 & 1 & 1 & \vline & 4
        \end{bmatrix}
        \rightarrow
        \begin{bmatrix}
            1 & 0 & 1 & \vline & 5 \\
            0 & 1 & -1 & \vline & 2 \\
            0 & 0 & 1 & \vline & 1
        \end{bmatrix}
        \rightarrow
            \begin{bmatrix}
            1 & 0 & 0 & \vline & 4 \\
            0 & 1 & 0 & \vline & 3 \\
            0 & 0 & 1 & \vline & 1
        \end{bmatrix}
        \implies
        \vec{[v]_\beta} = \begin{pmatrix}
            4 \\
            3 \\
            1
        \end{pmatrix}
    \]
}

\subsubsection{Find vektor ud fra koordinater med hensyn til base}
Givet \( \vec{[v]_\beta} = (c_1, c_2, \ldots, c_n)^T \), og \( \beta = \{ \vec{b_1}, \vec{b_2}, \ldots, \vec{b_n} \} \), findes \( \vec{v} \) ved:
\[
    \vec{v} = c_1 \cdot b_1 + c_2 \cdot b_2 + \ldots + c_n \cdot b_n
\]